% chapters/05/rangkuman.tex
\bagian{Rangkuman}

\begin{rangkumanbox}
\begin{itemize}
  \item \code{pygame.Vector2} menyatukan dua komponen posisi atau kecepatan
        menjadi satu nilai, sehingga penjumlahan dan perkalian dengan satu
        angka bekerja per komponen secara otomatis.
  \item Pada koordinat layar, sumbu tegak menuju ke bawah, sehingga gerak ke
        atas dinyatakan sebagai pengurangan nilai tegak.
  \item Keadaan tombol dibaca dengan \code{pygame.key.get_pressed()}, yang
        menghasilkan gerak berlangsung terus selama tombol ditahan; peristiwa
        \code{KEYDOWN} hanya terjadi sekali dan tidak memadai untuk keperluan
        itu.
  \item Vektor arah dinormalkan menjadi panjang satu sebelum dikalikan
        kecepatan. Tanpa langkah itu, gerak diagonal menjadi akar dua kali
        lebih cepat, yaitu 41 persen.
  \item \code{normalize()} mengembalikan vektor baru, sedangkan
        \code{normalize_ip()} mengubah vektor asalnya.
  \item Perpindahan dihitung sebagai kecepatan dikalikan selang waktu:
        \code{position = position + velocity * delta_seconds}. Inilah yang
        membuat kecepatan objek tidak bergantung pada laju bingkai.
  \item Urutan perhitungannya mengikuti kaidah Euler: kecepatan diperbarui
        lebih dahulu, baru posisi.
  \item Pembuktian melalui pengukuran memperlihatkan kecepatan terukur
        sebesar 200,0 piksel per detik pada laju 30, 60, dan 120 bingkai per
        detik.
  \item Jarak tempuh pada selang waktu yang sama dapat berbeda sedikit, karena
        \code{clock.tick()} hanya menjamin laju tidak melampaui batas, bukan
        bahwa setiap bingkai berlangsung tepat sepanjang jatahnya.
  \item Posisi yang sebenarnya disimpan sebagai pecahan, sedangkan
        \code{Rect} dipakai untuk menggambar dan memeriksa tabrakan. Menyimpan
        posisi sebagai bilangan bulat membuat objek tidak bergerak sama sekali
        pada kecepatan rendah, karena nilai pecahannya terpotong setiap
        bingkai.
\end{itemize}
\end{rangkumanbox}
